\chapter{System Design}

\section{System Architecture}

\project{} is a three-surface system with one source of truth. The phone is an
interface, the ESP32 is an actuator with a radio, and the backend owns identity,
parking state and money.

Three properties of this shape explain most of the code:

\begin{enumerate}
  \item \textbf{The application never authorises a gate.} It carries a session
        identifier over BLE and displays the answer; it holds no credential a gate
        would accept.
  \item \textbf{The gate never decides.} It holds a device secret that lets it ask,
        and a servo it may drive only after the answer arrives. Losing Wi-Fi closes
        gates rather than opening them.
  \item \textbf{The backend never touches hardware.} It returns a verdict and
        records the state change, which keeps the mechanism replaceable.
\end{enumerate}

\begin{table}[htbp]
  \centering
  \small
  \begin{tabularx}{\textwidth}{@{}lX@{}}
    \toprule
    \textbf{Component} & \textbf{Responsibility} \\
    \midrule
    Mobile application & Registration, login and token persistence; vehicle
      management; BLE discovery; session start; live fee; payment; gate commands;
      history. \\
    Backend & Identity and JWT, vehicles, facilities and pricing, session lifecycle,
      fee computation, mock payment, gate verdicts, device registry, remote command
      queue, stale-session sweep, health endpoint. \\
    PostgreSQL & Durable state for users, vehicles, facilities, sessions, payments,
      gates and commands; pricing stored as \code{jsonb}. \\
    Gate firmware & BLE peripheral and GATT service, Wi-Fi client, servo actuation,
      fixed-timer close, stall watchdog, remote-command poller. \\
    Serial bridge & Demonstration-only relay from the firmware's framed serial
      requests to the backend. \\
    Test workflow & Compose profile that boots dependencies, proves the production
      image is healthy, runs the suite and exports reports. \\
    \bottomrule
  \end{tabularx}
  \caption{Component responsibilities.}
  \label{tab:components}
\end{table}

\begin{table}[htbp]
  \centering
  \small
  \begin{tabularx}{\textwidth}{@{}lX@{}}
    \toprule
    \textbf{Decision} & \textbf{Rationale and consequence} \\
    \midrule
    BLE advertisement as the facility identifier & No GPS, QR or manual entry; the
      gate broadcasts \code{DWARA-<facilityCode>}. \\
    One gate serves entry and exit & The direction is a software state, not a second
      piece of hardware; the facility needs one servo. \\
    Backend authorises, firmware obeys & A forged client cannot open a gate; the
      firmware has no rule engine at all. \\
    Backend clock is the only clock & Neither the phone nor the ESP32 can move a
      bill; all timestamps are server-side. \\
    Per-second flat tariff & Trivial to explain and audit; stored per facility as
      JSON so prices change without a migration. \\
    One active session per user and per vehicle & Removes a class of abuse and
      simplifies the dashboard to a single active session. \\
    Payment freezes the fee & The exit gate asks about payment, not about elapsed
      time, so a queue at the exit cannot increase a bill. \\
    Fixed-timer gate close & No sensors; eight seconds is long enough for a car and
      short enough to limit tailgating. \\
    Abstract actuator and transport & Servo, LED, Wi-Fi and serial are
      interchangeable behind interfaces, which is what makes bench testing possible. \\
    Stale-session sweep & A session abandoned before the gate would otherwise block
      the account forever. \\
    \bottomrule
  \end{tabularx}
  \caption{Key design decisions.}
  \label{tab:key-decisions}
\end{table}

Deployment is a single Compose project. The development stack publishes two
services; the test stack adds a third behind a profile, so ordinary development is
unaffected.

\section{Use Cases}

\begin{table}[htbp]
  \centering
  \small
  \begin{tabularx}{\textwidth}{@{}lXl@{}}
    \toprule
    \textbf{ID} & \textbf{Use case} & \textbf{Actor} \\
    \midrule
    UC1 & Register and sign in & Driver \\
    UC2 & Manage vehicles & Driver \\
    UC3 & Discover a facility over BLE & Driver \\
    UC4 & Start a parking session & Driver \\
    UC5 & Open the entry gate & Driver, gate \\
    UC6 & View the running fee & Driver \\
    UC7 & Pay and open the exit gate & Driver, gate \\
    UC8 & Review session history & Driver \\
    UC9 & Provision a facility and a gate & Operator \\
    UC10 & Override a gate remotely & Operator, gate \\
    \bottomrule
  \end{tabularx}
  \caption{Use-case inventory.}
  \label{tab:use-cases}
\end{table}

The two use cases that differ most from a conventional system are UC5 and UC10. In
UC5 the driver's tap is a request, not an authorisation: the gate validates it with
the backend. In UC10 the operator never contacts the gate; the command is queued and
consumed by the gate's next poll.

\subsection*{Detailed use case: open the exit gate (UC7)}

\begin{table}[htbp]
  \centering
  \small
  \begin{tabularx}{\textwidth}{@{}lX@{}}
    \toprule
    \textbf{Field} & \textbf{Description} \\
    \midrule
    Precondition & The caller owns a session in \code{PAYMENT\_PENDING} with a
      successful payment; the phone is in the foreground near the gate. \\
    Main flow & 1. Driver taps \emph{Open exit gate}. 2. Application writes
      \code{\{sessionId, intent:"EXIT"\}} to the gate. 3. Gate calls
      \code{POST /gate/validate} with its device secret. 4. Backend verifies the
      device, facility, state and payment, stamps the exit and returns
      \code{open}. 5. Gate sweeps the servo, reports \code{OPENED}, holds eight
      seconds, closes. 6. Application refreshes the session and shows completion. \\
    Alternate flow & If the session is unpaid the backend returns
      \code{keep\_closed/PAYMENT\_REQUIRED}; the gate reports \code{DENIED} and the
      screen shows the payment prompt. \\
    Alternate flow & If the gate already exited within the grace window, the
      backend re-opens without changing state. \\
    Failure & Backend unreachable: the gate reports
      \code{DENIED/BACKEND\_UNREACHABLE} and stays closed. \\
    Postcondition & Session is \code{EXITED} (or unchanged for a grace re-open); the
      fee is frozen. \\
    \bottomrule
  \end{tabularx}
  \caption{Use-case description for the exit flow.}
  \label{tab:uc7}
\end{table}

\section{Layered Structure}

The backend is layered rather than organised as a single inheritance tree:
controllers depend on services, services on repositories and domain objects, and
every external capability is an interface.

\begin{table}[htbp]
  \centering
  \small
  \begin{tabularx}{\textwidth}{@{}lX@{}}
    \toprule
    \textbf{Package} & \textbf{Contents} \\
    \midrule
    \code{controllers} & Six REST controllers (\code{Auth}, \code{Vehicle},
      \code{Facility}, \code{Session}, \code{Gate}, \code{Device}) and one
      exception advice. \\
    \code{services} & Eleven service classes: identity, vehicles, facilities,
      sessions, billing, gate authorisation, provisioning, device authentication,
      commands, polling client, scheduler. Four interfaces. \\
    \code{repositories} & Seven Spring Data interfaces, including a locked user
      lookup, a locked session lookup and a paged history query. \\
    \code{entities} & Seven entities, eleven enums and the \code{PricingPolicy}
      record. \\
    \code{dtos} & Request records, response DTOs and the generic
      \code{ResponseObject} envelope. \\
    \code{configs} & Filter chain, JWT filter, admin-key filter and session timing
      properties. \\
    \bottomrule
  \end{tabularx}
  \caption{Layer contents.}
  \label{tab:layers}
\end{table}

\section{Session Flow}

\subsection{Entry}

The entry path has three decisions that are easy to get wrong. The session is created
\emph{before} the driver reaches the gate, so a failure at the gate does not lose the
vehicle; the gate validates the device's facility before reading the session; and the
entry timestamp is written inside the same transaction that decides \code{open}.

\subsection{Session state machine}

The session is the central object; every other rule refers to its state.

\begin{table}[htbp]
  \centering
  \small
  \begin{tabularx}{\textwidth}{@{}llX@{}}
    \toprule
    \textbf{From} & \textbf{To} & \textbf{Trigger and rule} \\
    \midrule
    (new) & \code{CREATED} & start session, active facility, no active session for the
      user or the vehicle \\
    \code{CREATED} & \code{PARKED} & entry verdict \code{open}; entry time stamped \\
    \code{CREATED} & \code{ABORTED} & explicit cancel, or the sweep after fifteen
      minutes \\
    \code{PARKED} & \code{PAYMENT\_PENDING} & payment succeeds; amount frozen \\
    \code{PAYMENT\_PENDING} & \code{EXITED} & exit verdict \code{open} with a
      successful payment \\
    \code{EXITED} & (re-open) & exit verdict inside the grace window; no state
      change \\
    \bottomrule
  \end{tabularx}
  \caption{Legal session transitions. Everything else is refused with a reason.}
  \label{tab:transitions}
\end{table}

\subsection{Gate decision table}

\begin{table}[htbp]
  \centering
  \small
  \begin{tabularx}{\textwidth}{@{}llX@{}}
    \toprule
    \textbf{Intent} & \textbf{State} & \textbf{Verdict} \\
    \midrule
    \code{ENTRY} & \code{CREATED} at the same facility & \code{open}, session becomes
      \code{PARKED} \\
    \code{ENTRY} & \code{PARKED}, \code{PAYMENT\_PENDING}, \code{EXITED} &
      \code{keep\_closed} / \code{WRONG\_STATE\_FOR\_ENTRY} \\
    \code{ENTRY} & \code{ABORTED} & \code{keep\_closed} / \code{STALE\_SESSION} \\
    \code{EXIT} & \code{PAYMENT\_PENDING} with a successful payment & \code{open},
      session becomes \code{EXITED} \\
    \code{EXIT} & \code{PAYMENT\_PENDING} without payment & \code{keep\_closed} /
      \code{PAYMENT\_REQUIRED} \\
    \code{EXIT} & \code{CREATED} or \code{PARKED} & \code{keep\_closed} /
      \code{PAYMENT\_REQUIRED} \\
    \code{EXIT} & \code{EXITED} inside the grace window & \code{open}, no state
      change \\
    \code{EXIT} & \code{EXITED} outside the grace window & \code{keep\_closed} /
      \code{STALE\_SESSION} \\
    any & unknown session & \code{keep\_closed} / \code{UNKNOWN\_SESSION} \\
    any & device at another facility & \code{keep\_closed} / \code{WRONG\_FACILITY} \\
    \bottomrule
  \end{tabularx}
  \caption{Decision table implemented by the gate service.}
  \label{tab:decisions}
\end{table}

\subsection{BLE interaction design}

The BLE link is a request channel: it carries a session identifier and an intent, and
returns the gate's answer. It carries no credential and no verdict the gate would
trust on its own. The service and characteristic layout follows the Generic
Attribute Profile of the Bluetooth Core Specification \cite{ble_gatt}.

\begin{table}[htbp]
  \centering
  \small
  \begin{tabularx}{\textwidth}{@{}lX@{}}
    \toprule
    \textbf{Element} & \textbf{Design} \\
    \midrule
    Advertisement & local name \code{DWARA-<facilityCode>}; service UUID
      \code{4fafc201-1fb5-459e-8fcc-c5c9c331914b} \\
    Command characteristic & \code{beb5483e-36e1-4688-b7f5-ea07361b26a8}, write \\
    Status characteristic & \code{c7d4a2d0-5f3a-4c3e-9b1a-8e2f6d4c5b7a}, read and
      notify \\
    Command payload & \code{\{"sessionId":"<uuid>","intent":"ENTRY"|"EXIT"|"CLOSE"\}},
      at most 128 bytes \\
    Status payload & \code{\{"sessionId":"<uuid>","result":"OPENED","reason":"NONE"\}} \\
    Timing & 800\,ms servo settle, 8\,s open hold, 3\,s stall watchdog, 5\,s command
      poll, 15\,s application round-trip budget \\
    \bottomrule
  \end{tabularx}
  \caption{BLE protocol design.}
  \label{tab:ble}
\end{table}

\section{Database Design}

Seven entities describe the domain. Primary keys are server-generated UUIDs, money is
decimal, and enumerations are stored as strings so a database inspection is readable
without a lookup table.

\begin{table}[htbp]
  \centering
  \small
  \begin{tabularx}{\textwidth}{@{}lX@{}}
    \toprule
    \textbf{Entity} & \textbf{Fields and role} \\
    \midrule
    \code{User} & \code{userId}, unique \code{email}, \code{name},
      \code{passwordHash}, \code{userRole}, \code{createdAt}. Table name
      \code{users}, because \code{user} is reserved in PostgreSQL. \\
    \code{Vehicle} & \code{vehicleId}, \code{numberPlate}, \code{vehicleType},
      owning user. \\
    \code{ParkingFacility} & \code{facilityId}, unique \code{facilityCode},
      \code{name}, \code{address}, \code{status}, \code{PricingPolicy} as
      \code{jsonb}. \\
    \code{ParkingSession} & \code{sessionId}, facility, user, vehicle,
      \code{sessionStatus}, \code{createdAt}, \code{entryTimestamp},
      \code{exitTimestamp}, \code{exitGateOpenedAt}, \code{totalFee}. \\
    \code{Payment} & \code{paymentId}, session, \code{amount}, \code{status},
      \code{transactionReference}. \\
    \code{GateController} & \code{controllerId}, facility, \code{gateType},
      \code{firmwareVersion}, write-only \code{deviceSecretHash}. \\
    \code{DeviceCommand} & \code{commandId}, device, \code{command}, \code{status},
      \code{createdAt}, \code{deliveredAt}. \\
    \bottomrule
  \end{tabularx}
  \caption{Entity design.}
  \label{tab:entities}
\end{table}

\begin{table}[htbp]
  \centering
  \small
  \begin{tabularx}{\textwidth}{@{}lX@{}}
    \toprule
    \textbf{Enumeration} & \textbf{Values and meaning} \\
    \midrule
    \code{SessionStatus} & \code{CREATED}, \code{PARKED}, \code{PAYMENT\_PENDING},
      \code{EXITED}, \code{ABORTED}; the first three form the active set. \\
    \code{PaymentStatus} & \code{SUCCESS}, with room for failure states. \\
    \code{VehicleType} & \code{MC}, \code{LMV}, \code{HMV}. \\
    \code{GateType} & \code{ENTRY}, \code{EXIT}. \\
    \code{GateIntent} & \code{ENTRY}, \code{EXIT}. \\
    \code{GateDecision} & \code{OPEN}, \code{KEEP\_CLOSED}. \\
    \code{DenyReason} & \code{NONE}, \code{STALE\_SESSION},
      \code{WRONG\_STATE\_FOR\_ENTRY}, \code{PAYMENT\_REQUIRED},
      \code{WRONG\_FACILITY}, \code{UNKNOWN\_SESSION}. \\
    \code{ParkingFacilityStatus} & \code{ACTIVE}, \code{INACTIVE},
      \code{MAINTENANCE}. \\
    \code{UserRole} & \code{USER}. \\
    \code{DeviceCommandType} & \code{NONE}, \code{OPEN}, \code{KEEP\_CLOSED}. \\
    \code{CommandStatus} & \code{PENDING}, \code{DELIVERED}. \\
    \bottomrule
  \end{tabularx}
  \caption{Enumerations.}
  \label{tab:enums}
\end{table}

Three modelling details carry most of the lifecycle safety:

\begin{itemize}
  \item \code{createdAt} is written when the user taps \emph{start}. It is what the
        stale-session sweep measures, because a session that never reached a gate has
        no entry time.
  \item \code{entryTimestamp} is written only by the entry verdict. It is the only
        input to billing, so a session with no entry is provably free.
  \item \code{exitGateOpenedAt} is separate from \code{exitTimestamp} in intent,
        though in practice they hold the same instant. The grace window is measured
        from the gate opening, not from payment, so paying early and exiting late
        does not forfeit the re-open.
\end{itemize}

\section{UI Design}

The application is deliberately thin: it renders state, collects intent and carries
messages. It does not decide fees or gate outcomes.

\begin{table}[htbp]
  \centering
  \small
  \begin{tabularx}{\textwidth}{@{}lX@{}}
    \toprule
    \textbf{Screen} & \textbf{Purpose and backend calls} \\
    \midrule
    Welcome & static onboarding \\
    Register & \code{POST /auth/register}, then login \\
    Login & \code{POST /auth/login}, then \code{GET /auth/me} \\
    Dashboard & active session and vehicles (\code{/sessions/active}, \code{/vehicles}) \\
    Discover & BLE scan, then \code{GET /facilities/\{code\}} \\
    Start Session & vehicle picker and \code{POST /sessions} \\
    Active Session & entry/exit BLE commands, live status \\
    Payment & fee polling and \code{POST /sessions/\{id\}/pay} \\
    Activity & \code{GET /sessions?page=0\&size=50} \\
    Profile & vehicle CRUD, sign-out, bench entry point \\
    Gate Test & manual BLE commands for diagnostics \\
    \bottomrule
  \end{tabularx}
  \caption{Screen design.}
  \label{tab:screens}
\end{table}

Navigation is a bottom tab bar over a native stack: the session screens are a stack
because they depend on each other, while discovery, activity and profile are
independent tabs. The payment screen polls the fee every two seconds while it is
visible and stops on unmount, so a background screen never accumulates requests.

\section{Security Design}

Three kinds of caller exist, and each has its own credential. They are not
interchangeable, and each is deliberately narrow.

\begin{table}[htbp]
  \centering
  \small
  \begin{tabularx}{\textwidth}{@{}lXl@{}}
    \toprule
    \textbf{Scheme} & \textbf{Used for} & \textbf{Storage} \\
    \midrule
    JWT bearer & Every user-facing endpoint. The subject is the user id; the token
      carries no role and no profile. & Stateless; signed with a base64 secret from
      the environment. \\
    \code{X-Admin-Key} & Facility and device administration. Compared before the
      authorization rules run. & Environment variable; blank closes the surface. \\
    Device bearer secret & \code{/gate/validate} and the command poll. 256 bits of
      entropy, returned once at registration. & BCrypt hash only. \\
    \bottomrule
  \end{tabularx}
  \caption{Credential schemes.}
  \label{tab:credentials}
\end{table}

The filter chain checks administration before routing and extracts JWTs before the
web authorization rules \cite{springsecurity}. A token is parsed only if it has three
dot-separated segments, because device secrets arrive in the same
\code{Authorization} header and must not be mistaken for tokens. A malformed or
expired token leaves the request anonymous rather than failing inside the filter; the
authorization rule then decides. Passwords and device secrets use the same one-way
treatment, so a database dump yields neither a usable password nor a usable gate
credential.

The token design follows the JSON Web Token standard \cite{rfc7519}: a signed
subject, an issue time and an expiry, with no sensitive claim in the payload. The
threat model below follows the structure recommended for API security assessments
\cite{owasp2023} and records accepted risk rather than hiding it.

\begin{table}[htbp]
  \centering
  \small
  \begin{tabularx}{\textwidth}{@{}>{\raggedright\arraybackslash}p{5.4cm}>{\raggedright\arraybackslash}X@{}}
    \toprule
    \textbf{Threat} & \textbf{Mitigation or accepted risk (v1)} \\
    \midrule
    Rogue BLE client writes another user's session id & Gate re-validates with the
      backend; worst case is \code{keep\_closed}. \\
    Replay at the exit gate & Exit requires payment and the right state;
      \code{EXITED} is terminal apart from one bounded grace re-open. \\
    Stolen JWT & Bearer-token risk bounded by the configured expiry; the client
      clears the token on 401/403. \\
    Device secret extracted from a flashed board & Only the hash is stored, but
      physical access to a provisioned unit is accepted as game over. \\
    BLE spoofing or man-in-the-middle & Accepted in version one; the backend is the
      trust anchor, not the link. \\
    Fee manipulation & All arithmetic is server-side and uses only the backend
      clock. \\
    Unauthenticated administration & Static key header, closed by default when the
      key is unset. \\
    Payment fraud & The provider is a mock; no money moves in version one. \\
    Tailgating & Fixed-timer close with no sensors; accepted and documented. \\
    \bottomrule
  \end{tabularx}
  \caption{Threat model for the version-one posture.}
  \label{tab:threats}
\end{table}
